% Part: proof-theory
% Chapter: cut-elimination
% Section: interpolation
%READERNOTE{SOL6-A296: चर-साक्षीचा प्रसंग भाषा-सीमेवरून पूर्ण केला आहे. सामान्य अंतर्वेशक सूत्ररूपात असतो; बंद संदर्भांतील उपयोगांसाठी मुक्त चरे सार्वत्रिकपणे बांधून वाक्यरूप घेतले आहे. OLINC-344 मधील ऐतिहासिक missing-case नोंद मूळ स्रोतासाठी कायम आहे, सध्याचा प्रसंग दुरुस्त आहे.}
%READERNOTE{SOL6-A303: बेथच्या परिभाष्यता परिच्छेदात उरलेल्या इंग्रजी predicate ऐवजी मराठी विधेय-संज्ञा वापरली आहे.}

\documentclass[../../../include/open-logic-section]{subfiles}

\begin{document}

\olfileid{pt}{cut}{itp}

\olsection{माएहाराचे उपप्रमेय आणि क्रेगचे अंतर्वेशन प्रमेय}

विल्यम क्रेग यांचे अंतर्वेशन प्रमेय असे सांगते:
$!A \Entails !B$ असेल, तर असे !!a{formula}~$!C$
असते (त्याला \emph{अंतर्वेशक} म्हणतात) की
$!A \Entails !C$ आणि $!C \Entails !B$.
विशेष म्हणजे, $!C$ मधील सर्व !!{constant}s आणि
!!{predicate}s हे $!A$ व $!B$ या दोन्हींमध्ये येतील
असे $!C$ निवडता येते. (येथे कोणतीही !!{function}s
नाहीत असे गृहीत धरतो.) ही अट नसेल, तर
$!A$ आणि $!B$ हेच सहजपणे अंतर्वेशक ठरतील.

अंतर्वेशन प्रमेय अभिजात प्रथम-क्रम तर्कशास्त्राला लागू होते.
त्याचा शोध प्रथम-क्रम तर्कशास्त्राचा “शेवटचा महत्त्वाचा गुणधर्म”
शोधला गेल्याप्रमाणे वर्णिला गेला आहे.
प्रतिमान उपपत्ती आणि स्वयंचलित तर्कपद्धती यांत त्याचे
महत्त्वाचे उपयोग आहेत. त्यामागील मूळ प्रेरणा आणि
उपयोग विज्ञानाच्या तत्त्वज्ञानात होता: एखाद्या उपपत्तीचे
स्वयंसिद्धकीकरण कोणत्या मूलभूत संज्ञांनी करता येते
किंवा येत नाही, हे तपासण्यासाठी ते वापरता येते.
यावरून बेथचे परिभाष्यता प्रमेयही निष्पन्न होते:
एखादी उपपत्ती संकल्पना अंतर्निहितपणे परिभाषित करू
शकत असेल, तर ती स्पष्टपणेही परिभाषित करू शकते.

गणिती तर्कशास्त्रातील अनेक निष्कर्षांप्रमाणे,
अंतर्वेशन प्रमेयाची सिद्धता प्रतिमान उपपत्तीच्या
तसेच सिद्धता उपपत्तीच्या पद्धतींनी करता येते.
सिद्धता-उपपत्तीय पद्धतीचा लाभ असा की ती
केवळ अंतर्वेशकाचे अस्तित्व दाखवत नाही, तर
$!A \Entails !B$ याच्या !!a{proof} पासून
अंतर्वेशक मोजून देणारी अल्गोरिदमही देते;
उदाहरणार्थ, $!A \Sequent !B$ याची~\Log{G3c}
मधील !!a{proof} वापरता येते.

$!A$ मधील !!{constant}s आणि !!{predicate}s
यांचा संच $\Lang L(!A)$ असू द्या आणि
$\Lang L(\Gamma) = \bigcup \Setabs{\Lang L(!A)}{!A \in
\Gamma}$ असू द्या. या संपूर्ण विभागात
!!{function}s नसलेली !!{formula}s आपण विचारात घेतो.

\begin{thm}[क्रेगचे अंतर्वेशन प्रमेय]\ollabel{thm:interpolation}
$!A$ ची !!{language} $\Lang L(!A)$ आणि
$!B$ ची $\Lang L(!B)$ असू द्या.
$\Log{G3c} \Proves !A \Sequent !B$ असेल, तर
$\Lang L(!C) \subseteq \Lang L(!A) \cap \Lang L(!B)$
या !!{language} मधील असे !!a{formula}~$!C$ असते की
$\Log{G3c} \Proves !A \Sequent !C$
आणि $\Log{G3c} \Proves !C \Sequent !B$.
\end{thm}

~\Log{G3c} मधील !!{proof}s च्या रचनेचा
उपयोग करण्यासाठी आपण अंतर्वेशन प्रमेयाचे
अधिक व्यापक विधान सिद्ध करू.
त्यासाठी क्रमवर्तींचे पूर्वांग व उत्तरांग
प्रत्येकी दोन भागांत विभागलेले मानू.
$\maeh{\Gamma_1}{\Gamma_2} \Sequent
\maeh{\Delta_1}{\Delta_2}$ या लेखनात
$\Gamma_1$ आणि $\Delta_1$ पहिल्या भागात,
तर $\Gamma_2$ आणि $\Delta_2$ दुसऱ्या
भागात आहेत. पुढील उपप्रमेयावरून
अंतर्वेशन प्रमेय लगेच मिळते.

\begin{lem}[माएहाराचे उपप्रमेय]\ollabel{lem:maehara}
$\Log{G3c}
\Proves \maeh{\Gamma_1}{\Gamma_2} \Sequent \maeh{\Delta_1}{\Delta_2}$
असेल, तर $\Lang L(!C)
\subseteq \Lang L(\Gamma_1,
\Delta_1) \cap \Lang L(\Gamma_2, \Delta_2)$
या !!{language} मधील असे !!a{formula}~$!C$ असते की
$\Log{G3c}
\Proves \Gamma_1 \Sequent \Delta_1, !C$ आणि
$\Log{G3c} \Proves !C,
\Gamma_2 \Sequent \Delta_2$.
\end{lem}
चारही संदर्भांतील सूत्रे वाक्ये असतील, तर अंतर्वेशकातील
मुक्त चरे सार्वत्रिकपणे बांधून वाक्यरूप अंतर्वेशक घेता येतो.
एका मुक्त चरासाठी पहिल्या सिद्धतेत तो चर आधी नव्या
स्थिरांकाने बदलून $\RightR{\lforall}$ लावायचा; बंद संदर्भांमुळे
आयगेन चराची अट पूर्ण होते. दुसऱ्या सिद्धतेत त्या चराचे
पद वापरून $\LeftR{\lforall}$ आणि आवश्यक दुर्बलीकरण लावायचे.
प्रत्येक मुक्त चरासाठी हे केल्याने दोन्ही क्रमवर्ती सिद्धतायोग्य
राहतात आणि स्थिरांक-विधेयांची भाषा-सीमा बदलत नाही.

\begin{proof}
  $\pi \Proves \maeh{\Gamma_1}{\Gamma_2} \Sequent
  \maeh{\Delta_1}{\Delta_2}$ असे गृहीत धरा.
  $n=\pheight{\pi}$ यावर विगमन करून विधान सिद्ध करू.

  $n=0$ असेल, तर
  $\maeh{\Gamma_1}{\Gamma_2} \Sequent
  \maeh{\Delta_1}{\Delta_2}$ हे स्वयंसिद्धक आहे
  आणि एखादे अणुरूप !!{formula}~$!A$
  हे $\Gamma_1, \Gamma_2$ आणि
  $\Delta_1, \Delta_2$ या दोन्ही बाजूंना येते.
  डावीकडे आणि उजवीकडे $!A$ पहिल्या की
  दुसऱ्या भागात आहे यावर चार प्रसंग ठरतात.
  \begin{enumerate}
    \item $!A \in \Gamma_1$ आणि $!A \in \Delta_1$.
    $\Gamma_1 \Sequent \Delta_1, !C$ आणि
    $!C, \Gamma_2 \Sequent \Delta_2$ या दोन्ही
    !!{provable} असतील असे~$!C$ हवे.
    $!A$ दोन्ही बाजूंना असल्याने, $!C$ काहीही
    निवडले तरी पहिली क्रमवर्ती आधीच स्वयंसिद्धक आहे.
    $!C, \Gamma_2 \Sequent \Delta_2$ ही दुसरी
    क्रमवर्ती !!{provable} असण्याची खात्री
    करण्यासाठी $!C \ident \lfalse$ घ्यावे लागते.
    \item $!A \in \Gamma_1$ आणि $!A \in \Delta_2$.
    मग $\Gamma_1 \Sequent \Delta_1, !A$ आणि
    $!A, \Gamma_2 \Sequent \Delta_2$ ही दोन्ही
    स्वयंसिद्धके आहेत; म्हणून $!C \ident !A$ घेता येते.
    \item $!A \in \Gamma_2$ आणि $!A \in \Delta_1$.
    मग $!A, \Gamma_1 \Sequent \Delta_1$ आणि
    $\Gamma_2 \Sequent \Delta_2, !A$ ही स्वयंसिद्धके
    असतील; पण त्यांत~$!A$ चुकीच्या बाजूंना आहे.
    म्हणून $!A$ अंतर्वेशक नाही.
    तथापि $\RightR{\lnot}$ आणि $\LeftR{\lnot}$
    वापरून $\Gamma_1
    \Sequent \Delta_1, \lnot!A$ आणि
    $\lnot !A, \Gamma_2 \Sequent
    \Delta_2$ या क्रमवर्ती !!{prove} करता येतात.
    \item $!A \in \Gamma_2$ आणि $!A \in \Delta_2$.
    हा सरावासाठीचा प्रसंग आहे.
  \end{enumerate}
  $\lfalse \in \Gamma_1$ किंवा
  $\lfalse \in \Gamma_2$ यांमुळे
  $\maeh{\Gamma_1}{\Gamma_2} \Sequent
  \maeh{\Delta_1}{\Delta_2}$ स्वयंसिद्धक असण्याचे
  प्रसंगही सरावासाठी सोडले आहेत.

  $n>0$ असेल, तर $\pi$ च्या शेवटी
  तार्किक अनुमान आहे. कोणता नियम वापरला
  आणि मुख्य सूत्र कोणत्या भागात आहे, यानुसार
  प्रसंग वेगळे करावे लागतील. प्रत्येक प्रसंगात
  $\pheight{\pi_1} < n$ असलेल्या सर्व
  !!{proof}s~$\pi_1$ साठी, त्यांची अंतिम
  क्रमवर्ती कोणत्याही प्रकारे दोन भागांत
  विभागली असली तरी, उपप्रमेय खरे आहे हे
  विगमन गृहीतक वापरता येते. काही विशिष्ट
  प्रसंग पाहू.
  \begin{enumerate}
    \item $\pi$ च्या शेवटी $\RightR{\lnot}$
    असून मुख्य सूत्र~$\lnot
    !A$ आहे. $\lnot !A$ पहिल्या की
    दुसऱ्या भागात आहे, यानुसार $\pi$ चा
    शेवट पुढीलपैकी एक असतो:
    \[
    \AxiomC{}
    \RightLabel{$\pi_1$}
    \Deduce$\maeh{!A, \Gamma_1}{\Gamma_2} \fCenter
      \maeh{\Delta_1}{\Delta_2}$
    \RightLabel{\RightR{\lnot}}
    \UnaryInf$\maeh{\Gamma_1}{\Gamma_2} \fCenter
      \maeh{\Delta_1, \lnot !A}{\Delta_2}$
    \DisplayProof
    \]
    किंवा
    \[
    \AxiomC{}
    \RightLabel{$\pi_1$}
    \Deduce$\maeh{\Gamma_1}{!A, \Gamma_2} \fCenter
      \maeh{\Delta_1}{\Delta_2}$
    \RightLabel{\RightR{\lnot}}
    \UnaryInf$\maeh{\Gamma_1}{\Gamma_2} \fCenter
      \maeh{\Delta_1}{\Delta_2, \lnot !A}$
    \DisplayProof
    \]
    मुख्य !!{formula}~$\lnot !A$ पहिल्या
    भागात असेल, तर आधारक्रमवर्तीतील सक्रिय
    !!{formula}~$!A$ देखील आपण पहिल्या
    भागातच ठेवले आहे. ही निवड आपल्याला
    करता येते. सक्रिय सूत्र दुसऱ्या भागात
    ठेवले तरी $\pi_1$ च्या
    अंतिम क्रमवर्तीला विगमन गृहीतक लागू होईल;
    मात्र मग अंतर्वेशक शोधता येणार नाही
    (हे करून पाहा).

    आधी $\lnot !A$ पहिल्या भागात
    असलेला प्रसंग पाहू. विगमन गृहीतकानुसार
    $\pi_1$ च्या अंतिम क्रमवर्तीला
    अंतर्वेशक असतो, म्हणजे असे
    !!a{formula}~$!C_1$ असते की पुढील दोन्ही
    क्रमवर्ती
    \begin{align*}
    !A, \Gamma_1 & \Sequent \Delta_1, !C_1 &&\text{आणि} &
    !C_1, \Gamma_2 & \Sequent \Delta_2
    \intertext{सिद्ध करता येतात. आता आपल्याला असे !!a{formula}~$!C$ हवे की
    पुढील दोन्ही !!{provable} असतील:}
    \Gamma_1 & \Sequent \Delta_1, \lnot !A, !C &&\text{आणि} &
    !C, \Gamma_2 & \Sequent \Delta_2
    \end{align*}
    स्पष्टपणे $!C \ident !C_1$ घेता येते:
    उजवीकडील क्रमवर्ती !!{provable} आहे
    हे विगमन गृहीतकाने आधीच मिळते, आणि
    डावीकडील क्रमवर्ती विगमन गृहीतकातून
    मिळालेल्या क्रमवर्तीवर~$\RightR{\lnot}$
    लावून मिळते. विगमन गृहीतकाने
    $\Lang L(!C_1) \subseteq \Lang L(!A, \Gamma_1,
    \Delta_1) \cap \Lang L(\Gamma_2, \Delta_2)$
    हेदेखील मिळते. $\Lang L(\lnot
    !A) = \Lang L(!A)$ आणि
    $\Lang L(!C) = \Lang L(!C_1)$ असल्याने
    $\Lang L(!C_1) \subseteq \Lang L(\Gamma_1, \Delta_1, \lnot !A) \cap \Lang
    L(\Gamma_2, \Delta_2)$.

    दुसऱ्या प्रसंगात $\lnot !A$
    दुसऱ्या भागात असल्याने स्थिती थोडी
    वेगळी आहे. आधारक्रमवर्तीतील सक्रिय
    !!{formula} देखील दुसऱ्या भागात
    ठेवून विगमन गृहीतक लावतो; त्यातून
    \begin{align*}
    \Gamma_1 & \Sequent \Delta_1, !C_1 &&\text{आणि} &
    !C_1, !A, \Gamma_2 & \Sequent \Delta_2
    \intertext{सिद्ध करता येतात. यांतील दुसऱ्या क्रमवर्तीवर पुन्हा \RightR{\lnot} लावल्यास पुढील !!{proof}s मिळतात:}
    \Gamma_1 & \Sequent \Delta_1, !C_1 &&\text{आणि} &
    !C_1, \Gamma_2 & \Sequent \Delta_2, \lnot !A
    \end{align*}
    $!C_1$ अंतर्वेशक असेल तर नेमक्या
    याच क्रमवर्ती !!{provable} हव्या होत्या;
    म्हणून पुन्हा $!C \ident !C_1$ घेऊ.
    पूर्वीप्रमाणे $\Lang
    L(!C_1) \subseteq \Lang L(\Gamma_1, \Delta_1) \cap \Lang L(\Gamma_2,
    \Delta_2, \lnot !A)$.
    \item $\pi$ च्या शेवटी $\RightR{\lif}$
    असून मुख्य !!{formula}~$!A
    \lif !B$ आहे. $!A \lif
    !B$ पहिल्या की दुसऱ्या भागात आहे
    यावर पुन्हा दोन प्रसंग मिळतात.
    !!{proof}~$\pi$ चा शेवट असा असतो:
    \[
    \AxiomC{}
    \RightLabel{$\pi_1$}
    \Deduce$\maeh{!A, \Gamma_1}{\Gamma_2} \fCenter
      \maeh{\Delta_1, !B}{\Delta_2}$
    \RightLabel{\RightR{\lif}}
    \UnaryInf$\maeh{\Gamma_1}{\Gamma_2} \fCenter
      \maeh{\Delta_1, !A \lif !B}{\Delta_2}$
    \DisplayProof
    \]
    किंवा
    \[
    \AxiomC{}
    \RightLabel{$\pi_1$}
    \Deduce$\maeh{\Gamma_1}{!A, \Gamma_2} \fCenter
      \maeh{\Delta_1}{\Delta_2, !B}$
    \RightLabel{\RightR{\lif}}
    \UnaryInf$\maeh{\Gamma_1}{\Gamma_2} \fCenter
      \maeh{\Delta_1}{\Delta_2, !A \lif !B}$
    \DisplayProof
    \]
    येथेही आधारक्रमवर्तीतील सक्रिय
    !!{formula}s निष्कर्षातील मुख्य
    !!{formula} च्याच भागात ठेवली आहेत.
    पहिल्या प्रसंगात विगमन गृहीतकाने
    पुढील !!{proof}s मिळतात:
    \begin{align*}
    !A, \Gamma_1 & \Sequent \Delta_1, !B, !C_1 &&\text{आणि} & !C_1, \Gamma_2 & \Sequent \Delta_2
    \intertext{डावीकडील क्रमवर्ती \RightR{\lif} साठी योग्य आधारक्रमवर्ती आहे. म्हणून पुढील क्रमवर्ती !!{provable} आहेत:}
    \Gamma_1 & \Sequent \Delta_1, !A \lif !B, !C_1 &&\text{आणि} & !C_1, \Gamma_2 & \Sequent \Delta_2
    \end{align*}
    म्हणून $!C_1$ हा $\pi$ च्या
    अंतिम क्रमवर्तीचाही अंतर्वेशक आहे;
    पुन्हा $!C \ident !C_1$ घेता येते.

    दुसरा प्रसंगही याच प्रकारे हाताळता येतो.
    \item $\pi$ च्या शेवटी $\LeftR{\lif}$
    असून मुख्य !!{formula}~$!A
    \lif !B$ आहे. $!A \lif !B$
    पहिल्या भागात असेल, तर
    !!{proof}~$\pi$ चा शेवट असा असतो:
    \[
    \AxiomC{}
    \RightLabel{$\pi_1$}
    \Deduce$\maeh{\Gamma_1}{\Gamma_2} \fCenter
      \maeh{\Delta_1, !A}{\Delta_2}$
    \AxiomC{}
    \RightLabel{$\pi_2$}
    \Deduce$\maeh{!B, \Gamma_1}{\Gamma_2} \fCenter
      \maeh{\Delta_1}{\Delta_2}$
    \BinaryInf$\maeh{!A \lif !B, \Gamma_1}{\Gamma_2} \fCenter
      \maeh{\Delta_1}{\Delta_2}$
    \DisplayProof
    \]
    विगमन गृहीतकाने आता चार
    !!{provable} क्रमवर्ती मिळतात:
    डाव्या आधारक्रमवर्तीच्या
    अंतर्वेशक~$!C_1$ साठी दोन आणि
    दुसऱ्या आधारक्रमवर्तीच्या
    अंतर्वेशक~$!C_2$ साठी दोन:
    \begin{align*}
    \Gamma_1 & \Sequent \Delta_1, !A, !C_1 &&\text{(1)} &
    !C_1, \Gamma_2 & \Sequent \Delta_2 && \text{(2)}\\
    !B, \Gamma_1 & \Sequent \Delta_1, !C_2 &&\text{(3)} &
    !C_2, \Gamma_2 & \Sequent \Delta_2 && \text{(4)}
    \end{align*}
    दुर्बलीकरण लावून (1) च्या
    उजवीकडे~$!C_2$ आणि (3) च्या
    उजवीकडे~$!C_1$ जोडल्यास
    (1) आणि (3) या क्रमवर्ती
    \LeftR{\lif} च्या आधारक्रमवर्ती
    म्हणून वापरता येतात:
    \[
    \AxiomC{}
    \Deduce$\Gamma_1 \fCenter \Delta_1, !A, !C_1, !C_2$
    \AxiomC{}
    \Deduce$!B, \Gamma_1 \fCenter \Delta_1, !C_1, !C_2$
    \RightLabel{\LeftR{\lif}}
    \BinaryInf$!A \lif !B, \Gamma_1 \fCenter \Delta_1, !C_1, !C_2$
    \DisplayProof
    \]
    $\RightR{\lor}$ लावल्याने पुढील
    क्रमवर्ती मिळते:
    \[
    !A \lif !B, \Gamma_1 \fCenter \Delta_1, !C_1 \lor !C_2
    \]
    त्यामुळे या प्रसंगात
    $!C_1 \lor !C_2$ अंतर्वेशक असू
    शकतो असे सुचते. उरलेल्या
    (2) आणि~(4) क्रमवर्ती वापरून
    आवश्यक दुसरी क्रमवर्तीही
    !!{prove} करता येते:
    \[
    \AxiomC{}
    \Deduce$!C_1, \Gamma_2 \fCenter \Delta_2$
    \AxiomC{}
    \Deduce$!C_2, \Gamma_2 \fCenter \Delta_2$
    \RightLabel{\LeftR{\lor}}
    \BinaryInf$!C_1 \lor !C_2, \Gamma_2 \fCenter \Delta_2$
    \DisplayProof
    \]
    $!C_1 \lor !C_2$ योग्य
    भाषेत आहे हे तपासणे उरते.
    विगमन गृहीतकानुसार
    \begin{align*}
    \Lang L(!C_1) & \subseteq
    \Lang L(\Gamma_1, \Delta_1, !A) \cap \Lang L(\Gamma_2, \Delta_2) &\text{आणि}\\
    \Lang L(!C_2) & \subseteq
    \Lang L(!B, \Gamma_1, \Delta_1) \cap \Lang L(\Gamma_2, \Delta_2)
    \intertext{कारण}
    \Lang L(!A) & \subseteq \Lang L(!A \lif !B) &\text{आणि}\\
    \Lang L(!B) & \subseteq \Lang L(!A \lif !B)
    \intertext{म्हणून या दोन्ही अटी मिळतात}
    \Lang L(!C_1) & \subseteq
    \Lang L(\Gamma_1, \Delta_1, !A \lif !B) \cap \Lang L(\Gamma_2, \Delta_2) &\text{आणि}\\
    \Lang L(!C_2) & \subseteq
    \Lang L(\Gamma_1, \Delta_1, !A \lif !B) \cap \Lang L(\Gamma_2, \Delta_2)
    \intertext{आणि परिणामी, $\Lang L (!C_1 \lor !C_2) = {}$}
    \Lang L(!C_1) \cup \Lang L(!C_2) & \subseteq
    \Lang L(\Gamma_1, \Delta_1, !A \lif !B) \cap \Lang L(\Gamma_2, \Delta_2),
    \end{align*}
    आणि हेच अपेक्षित आहे.

    $!A \lif !B$ दुसऱ्या
    भागात असण्याचा प्रसंग
    वेगळा आहे, पण तशाच
    पद्धतीने हाताळता येतो.
    विगमन गृहीतकाने पुन्हा
    चार !!{provable} क्रमवर्ती मिळतात:
    \begin{align*}
    \Gamma_1 & \Sequent \Delta_1, !C_1 &&\text{(1)} &
    !C_1, \Gamma_2 & \Sequent \Delta_2, !A && \text{(2)}\\
    \Gamma_1 & \Sequent \Delta_1, !C_2 &&\text{(3)} &
    !C_2, !B, \Gamma_2 & \Sequent \Delta_2 && \text{(4)}
    \end{align*}
    आता (2) आणि (4) या क्रमवर्ती,
    काही !!{formula}s दुर्बल करून
    जोडल्यावर,~\LeftR{\lif}
    च्या आधारक्रमवर्ती ठरतात:
    \[
    \AxiomC{}
    \Deduce$!C_1, !C_2, \Gamma_2 \fCenter \Delta_2, !A$
    \AxiomC{}
    \Deduce$!B, !C_1, !C_2, \Gamma_2 \fCenter \Delta_2$
    \RightLabel{\LeftR{\lif}}
    \BinaryInf$!A \lif !B, !C_1, !C_2, \Gamma_2 \fCenter \Delta_2$
    \DisplayProof
    \]
    यांतून $!C_1$ आणि $!C_2$
    वापरून बनलेले एकच !!{formula}
    असणारी क्रमवर्ती हवी आहे;
    पण यावेळी ते पूर्वांगात हवे
    (कारण आता दुसऱ्या भागाशी
    संबंध आहे). $\LeftR{\land}$
    लावल्याने ते साधते; म्हणून
    $!C \ident
    (!C_1 \land !C_2)$ अंतर्वेशक घेऊ.
    उरलेल्या (1) आणि (3)
    क्रमवर्ती वापरून पहिल्या
    भागाची संबंधित क्रमवर्तीही
    !!{prove} करता येते:
    \[
    \AxiomC{}
    \Deduce$\Gamma_1 \fCenter \Delta_1, !C_1$
    \AxiomC{}
    \Deduce$\Gamma_1 \fCenter \Delta_1, !C_2$
    \RightLabel{\RightR{\land}}
    \BinaryInf$\Gamma_1 \fCenter \Delta_1, !C_1 \land !C_2$
    \DisplayProof
    \]
    पूर्वीप्रमाणे
    $\Lang L(!C_1 \land !C_2) \subseteq \Lang
    L(\Gamma_1, \Delta_1) \cap \Lang L(\Gamma_2, \Delta_2, !A \lif
    !B)$ मिळते.

    \item $\pi$ च्या शेवटी $\RightR{\lforall}$
    असून मुख्य !!{formula}~$\lforall[x][!A(x)]$ आहे:
    \[
    \AxiomC{}
    \RightLabel{$\pi_1$}
    \Deduce$\maeh{\Gamma_1}{\Gamma_2} \fCenter
    \maeh{\Delta_1, !A(c)}{\Delta_2}$
    \RightLabel{\RightR{\lforall}}
    \UnaryInf$\maeh{\Gamma_1}{\Gamma_2} \fCenter
      \maeh{\Delta_1, \lforall[x][!A(x)]}{\Delta_2}$
    \DisplayProof
    \]
    किंवा
    \[
    \AxiomC{}
    \RightLabel{$\pi_1$}
    \Deduce$\maeh{\Gamma_1}{\Gamma_2} \fCenter
      \maeh{\Delta_1}{\Delta_2, !A(c)}$
    \RightLabel{\RightR{\lforall}}
    \UnaryInf$\maeh{\Gamma_1}{\Gamma_2} \fCenter
      \maeh{\Delta_1}{\Delta_2, \lforall[x][!A(x)]}$
    \DisplayProof
    \]
    पहिल्या प्रसंगात विगमन गृहीतकाने
    $\Gamma_1 \Sequent \Delta_1, !A(c), !C_1$
    आणि $!C_1, \Gamma_2
    \Sequent \Delta_2$ यांच्या !!{proof}s मिळतात.
    पहिल्या क्रमवर्तीवर $\RightR{\lforall}$
    लावून $\Gamma_1 \Sequent \Delta_1, \lforall[x][!A(x)], !C_1$
    मिळते. ($\pi$ च्या अंतिम
    \RightR{\lforall} अनुमानात आयगेन चराची
    अट आधीच पूर्ण असल्याने येथेही ती पूर्ण आहे.)
    $\Lang
    L(!C_1) \subseteq \Lang L(\Gamma_1, \Delta_1, \lforall[x][!A(x)])
    \cap \Lang L(\Gamma_2, \Delta_2)$
    ही अट पूर्ण असेल, तर $!C_1$ अंतर्वेशक ठरेल.
    विगमन गृहीतकाने
    $\Lang L(!C_1) \subseteq \Lang L(\Gamma_1,
    \Delta_1, !A(c)) \cap \Lang L(\Gamma_2, \Delta_2)$
    मिळते. म्हणून~$c$ ची काळजी घ्यावी लागते:
    $!C_1$ मध्ये $c$ येऊ शकतो का? येऊ शकत नाही.
    $!C_1$ मध्ये $c$ आल्यास
    $c \in \Lang L(\Gamma_1, \Delta_1,
    !A(c))$ (जे खरे आहे) आणि
    $c \in \Lang L(\Gamma_2,
    \Delta_2)$ हेही दोन्ही खरे ठरतील.
    पण मग $\pi$ मधील
    $\RightR{\lforall}$ अनुमानाच्या
    निष्कर्षात $c$ येईल आणि
    आयगेन चराची अट मोडेल.

    दुसरा प्रसंगही यासारखाच आहे.

    \item $\pi$ च्या शेवटी
    $\RightR{\lexists}$ असून मुख्य
    !!{formula}~$\lexists[x][!A(x)]$ आहे.
    येथेही दोन प्रसंग आहेत.
    $\lexists[x][!A(x)]$ पहिल्या
    भागात असलेला प्रसंग पाहू;
    दुसरा सरावासाठी सोडला आहे.

    पहिल्या प्रसंगात $\pi$ चा
    शेवट असा असतो:
    \[
    \AxiomC{}
    \RightLabel{$\pi_1$}
    \Deduce$\maeh{\Gamma_1}{\Gamma_2} \fCenter
      \maeh{\Delta_1, \lexists[x][!A(x)], !A(t)}{\Delta_2}$
    \RightLabel{\RightR{\lexists}}
    \UnaryInf$\maeh{\Gamma_1}{\Gamma_2} \fCenter
      \maeh{\Delta_1, \lexists[x][!A(x)]}{\Delta_2}$
    \DisplayProof
    \]
    विगमन गृहीतकाने
    $\Gamma_1 \Sequent
    \Delta_1, \lexists[x][!A(x)], !A(t), !C_1$
    आणि $!C_1, \Gamma_2
    \Sequent \Delta_2$ यांच्या !!{proof}s मिळतात.
    पहिल्या क्रमवर्तीवर $\RightR{\lexists}$
    लावल्याने $\Gamma_1 \Sequent \Delta_1, \lexists[x][!A(x)], !C_1$
    मिळते. म्हणून
    $\Lang L(!C_1) \subseteq \Lang L(\Gamma_1, \Delta_1,
    \lexists[x][!A(x)]) \cap \Lang L(\Gamma_2, \Delta_2)$
    ही अट पूर्ण असल्यास $!C_1$ पुन्हा
    अंतर्वेशक ठरेल. येथे विगमन गृहीतकाने
    $\Lang L(!C_1)
    \subseteq \Lang L(\Gamma_1, \Delta_1, \lexists[x][!A(x)], !A(t))
    \cap \Lang L(\Gamma_2, \Delta_2)$
    मिळते. आपण !!{function}s नाहीत असे
    गृहीत धरले आहे, हे लक्षात घ्या.
    म्हणून पद~$t$ हे चर किंवा स्थिरांक असते.
    $t$ चर असेल, तर त्याच्या प्रतिस्थापनाने नवीन
    स्थिरांक किंवा विधेय-चिन्ह येत नाही:
    $\Lang L(!A(t))\subseteq\Lang L(!A(x))$.
    त्यामुळे वरील विगमन गृहीतकातील भाषा-सीमा
    अंतिम क्रमवर्तीच्या दोन भागांच्या भाषांच्या
    छेदातच राहते आणि $!C\ident !C_1$ घेता येते.
    $\RightR{\lexists}$ साठी आयगेन चराची अट नाही.
    आता $t$ हे !!a{constant} असल्याचा प्रसंग,
    म्हणजे $t\ident b$, पाहू.
    $b \notin \Lang L(!C_1)$ असेल,
    किंवा $b \in \Lang L(\Gamma_1, \Delta_1, \lexists[x][!A(x)]) \cap
    \Lang L(\Gamma_2, \Delta_2)$ असेल,
    तर अडचण नाही: $!C_1$ अंतर्वेशक घेता येते.
    पण $b \in \Lang
    L(!C_1)$ आणि $b \notin \Lang L(\Gamma_1, \Delta_1,
    \lexists[x][!A(x)]) \cap \Lang L(\Gamma_2, \Delta_2)$
    असेल तर? प्रथम~$!C_1$ मधील
    $b$ ची प्रत्येक उपस्थिती
    $!C_1$ मध्ये न येणारा नवा !!a{variable}~$y$
    निवडून त्याने बदलल्यावर
    $!C_1 \ident \Subst{!C_1'}{b}{y}$
    असे लिहिता येते; येथे $!C_1'$
    मध्ये~$b$ येत नाही. शिवाय,
    $b \notin
    \Lang L(\Gamma_1, \Delta_1, \lexists[x][!A(x)])$
    किंवा $b \notin
    \Lang L(\Gamma_2, \Delta_2)$.
    त्यानुसार $!C \ident
    \lforall[y][!C_1'(y)]$ किंवा
    $!C \ident \lexists[y][!C_1'(y)]$
    घेता येते. पहिल्या प्रसंगात:
    \[
    \AxiomC{}
    \Deduce$\Gamma_1 \fCenter
      \Delta_1, \lexists[x][!A(x)], !A(b), !C_1'(b)$
    \RightLabel{\RightR{\lexists}}
    \UnaryInf$\Gamma_1 \fCenter \Delta_1, \lexists[x][!A(x)], !C_1'(b)$
    \RightLabel{\RightR{\lforall}}
    \UnaryInf$\Gamma_1 \fCenter
      \Delta_1, \lexists[x][!A(x)], \lforall[y][!C_1'(y)]$
    \DisplayProof
    \]
    आणि
    \[
    \AxiomC{}
    \Deduce$!C_1'(b), \lforall[y][!C_1'(y)], \Gamma_2 \fCenter \Delta_2$
    \RightLabel{\LeftR{\lforall}}
    \UnaryInf$\lforall[y][!C_1'(y)], \Gamma_2 \fCenter \Delta_2$
    \DisplayProof
    \]
    सर्वात वरच्या क्रमवर्तींच्या
    !!{proof}s विगमन गृहीतकातून मिळतात
    (दुसऱ्या सिद्धतेत पूर्वांगात
    $\lforall[y][!C_1'(y)]$ जोडण्यासाठी
    दुर्बलीकरणाची ग्राह्यता वापरतो).
    पहिल्या !!{proof} मधील
    \RightR{\lforall} साठी आयगेन चराची
    अट पूर्ण आहे, कारण
    $b \notin \Lang L(\Gamma_1, \Delta_1,
    \lexists[x][!A(x)])$ आणि
    $!C_1'$ मध्ये~$b$ येत नाही,
    अशीच त्याची व्याख्या केली आहे.

    दुसऱ्या प्रसंगात, विगमन गृहीतक
    आणि $\lexists[y][!C_1'(y)]$
    जोडणाऱ्या दुर्बलीकरणाच्या
    ग्राह्यतेने पुढील मिळते:
    \[
    \AxiomC{}
    \Deduce$\Gamma_1 \fCenter
      \Delta_1, \lexists[x][!A(x)], !A(b), \lexists[y][!C_1'(y)], !C_1'(b)$
    \RightLabel{\RightR{\lexists}}
    \UnaryInf$\Gamma_1 \fCenter
      \Delta_1, \lexists[x][!A(x)], \lexists[y][!C_1'(y)], !C_1'(b)$
    \RightLabel{\RightR{\lexists}}
    \UnaryInf$\Gamma_1 \fCenter
      \Delta_1, \lexists[x][!A(x)], \lexists[y][!C_1'(y)]$
    \DisplayProof
    \]
    आणि
    \[
    \AxiomC{}
    \Deduce$!C_1'(b), \Gamma_2 \fCenter \Delta_2$
    \RightLabel{\LeftR{\lexists}}
    \UnaryInf$\lexists[y][!C_1'(y)], \Gamma_2 \fCenter \Delta_2$
    \DisplayProof
    \]
    दुसऱ्या !!{proof} मधील
    \LeftR{\lexists} साठी आयगेन चराची
    अट पूर्ण आहे, कारण
    $b \notin \Lang L(\Gamma_2,
    \Delta_2)$ आणि $!C_1'$ मध्ये~$b$
    येत नाही, अशीच त्याची व्याख्या केली आहे.
  \end{enumerate}
  उरलेले प्रसंगही यांसारखे आहेत;
  ते सरावासाठी सोडले आहेत.
\end{proof}

\begin{prob}
$\lnand$ हा संयोजक पुढील नियमांसह
आपल्याकडे आहे असे समजा:
\[
\Axiom$\Gamma \fCenter \Delta, !A$
\Axiom$\Gamma \fCenter \Delta, !B$
\RightLabel{\LeftR{\lnand}}
\BinaryInf$!A \lnand !B, \Gamma \fCenter \Delta$
\DisplayProof
\quad
\Axiom$!A, !B, \Gamma \fCenter \Delta$
\RightLabel{\RightR{\lnand}}
\UnaryInf$\Gamma \fCenter \Delta, !A \lnand !B$
\DisplayProof
\]
ज्या प्रसंगात $\pi$ चा शेवट
यांपैकी एखाद्या नियमाने होतो,
ते \olref{lem:maehara} च्या सिद्धतेत
पूर्ण करून माएहाराचे उपप्रमेय
अजूनही खरे आहे हे दाखवा.
\end{prob}

माएहाराचे उपप्रमेय सिद्ध झाल्यावर
त्यावरून क्रेगचे अंतर्वेशन प्रमेय
सिद्ध करता येते.

\begin{proof}[\olref{thm:interpolation} ची सिद्धता]
  $!A \Sequent !B$ !!{provable} आहे असे
  समजा. $\maeh{!A}{\ } \Sequent \maeh{\ }{!B}$
  याला \olref{lem:maehara} लावून
  अंतर्वेशक~$!C$ मिळवा. मग
  $\Proves !A \Sequent !C$ आणि
  $\Proves !C
  \Sequent !B$.
\end{proof}

$\Gamma$ ही !!{sentence}s चा संच असलेली
उपपत्ती, !!a{language}~$\Lang L$ मध्ये आहे
आणि त्यात $n$-स्थानी {विधेय}~$R$ आहे
असे समजा. पुढील अट पूर्ण असेल, तेव्हा
आणि केवळ तेव्हाच, $\Gamma$ ही
$R$ ला \emph{अंतर्निहितपणे परिभाषित करते}:
\begin{align*}
\Gamma, \Gamma'
& \Entails \lforall[x_1][\dots\lforall[x_n][(R(x_1,\dots,x_n) \liff
R'(x_1, \dots, x_n))]],
\intertext{येथे $\Gamma'$ म्हणजे $\Gamma$ मध्ये $R$ ची प्रत्येक उपस्थिती
$R' \notin \Lang L$ ने बदलल्यावर मिळणारी उपपत्ती. $\Gamma$ ही $R$ ला
\emph{स्पष्टपणे परिभाषित करते}, तेव्हा आणि केवळ तेव्हाच, $R$ नसलेले असे
$!A(x_1, \dots, x_n)$ असते की}
\Gamma &\Entails
\lforall[x_1][\dots\lforall[x_n][(R(x_1,\dots,x_n) \liff !A(x_1,
\dots, x_n))]].
\end{align*}
स्पष्टपणे, $\Gamma$ ही $R$ ला
स्पष्टपणे परिभाषित करत असेल, तर
ती $R$ ला अंतर्निहितपणेही परिभाषित करते.
उलट विधान लगेच स्पष्ट होत नसले,
तरी ते खरे आहे:

\begin{lem}[बेथचे परिभाष्यता प्रमेय]
  $\Gamma$ ही $R$ ला अंतर्निहितपणे
  परिभाषित करत असेल, तर ती त्याला
  स्पष्टपणेही परिभाषित करते.
\end{lem}

\begin{proof}
  $\Gamma$ ही $R$ ला अंतर्निहितपणे
  परिभाषित करते असे समजा. वरीलप्रमाणे
  $R'$ आणि $\Gamma'$ घ्या;
  $\Lang L' = \Lang L(\Gamma') = \Lang L \setminus \{R\}
  \cup \{R'\}$ असू द्या. गृहीतकानुसार
  \begin{align*}
    \Gamma, \Gamma' &
    \Sequent \lforall[x_1][\dots\lforall[x_n][(R(x_1,\dots,x_n)
    \liff R'(x_1, \dots, x_n))]]
  \intertext{$c_1$, \dots,~$c_n$ हे $\Gamma$ मध्ये न येणारे !!{constant}s असू द्या. व्युत्क्रमणीयतेच्या उपप्रमेयाने मिळते:}
    \Gamma, \Gamma', R(c_1,\dots,c_n) & \Sequent  R'(c_1, \dots, c_n)
  \intertext{माएहाराचे उपप्रमेय पुढील क्रमवर्तीला लावा:}
    \maeh{\Gamma, R(c_1,\dots,c_n)}{\Gamma'} & \Sequent
      \maeh{\ }{R'(c_1, \dots, c_n)}
  \intertext{या बंद संदर्भांसाठी आवश्यक मुक्त चरे सार्वत्रिकपणे बांधून वाक्यरूप अंतर्वेशक निवडा. मग नव्या स्थिरांकांच्या जागी संबंधित चरे ठेवून असे~$!A \ident !A(x_1, \dots, x_n)$ मिळते की}
    \Gamma, R(c_1,\dots,c_n) & \Sequent !A(c_1, \dots, c_n) \\
    !A(c_1, \dots, c_n), \Gamma' & \Sequent R'(c_1, \dots, c_n)
  \intertext{यांच्या !!{proof}s असतात. $\Lang L(!A) \subseteq
  (\Lang L(\Gamma) \setminus \{R\}) \cup \{c_1,\dots,c_n\}$ असल्याने $!A$ मध्ये~$R$ नसतो.
  दुसऱ्या क्रमवर्तीच्या !!{proof} मध्ये सर्वत्र $R'$ च्या जागी~$R$ ठेवून
  पुढील !!a{proof} मिळते:}
    !A(c_1, \dots, c_n), \Gamma & \Sequent R(c_1, \dots, c_n)
  \end{align*}
  येथील नव्या स्थिरांकांच्या जागी
  संबंधित चरे ठेवता येतात.
  \RightR{\lif} आणि
  \RightR{\lforall} लावून पुढील
  !!a{proof} मिळते:
  \[
    \Gamma \Sequent \lforall[x_1][\dots\lforall[x_n][(R(x_1,\dots,x_n) \liff !A(x_1, \dots, x_n))]].
  \]
  म्हणून $!A(x_1, \dots, x_n)$ हे
  $\Gamma$ मध्ये $!R$ ला स्पष्टपणे
  परिभाषित करते.
\end{proof}

अंतर्वेशनाशी समतुल्य असलेला
आणि त्याप्रमाणेच माएहाराचे उपप्रमेय वापरून
सिद्ध करता येणारा तिसरा निष्कर्ष असा:
दोन सुसंगत उपपत्ती~$\Gamma_1$ आणि
$\Gamma_2$ यांत $\Lang L(\Gamma) = \Lang L(\Gamma_1)
\cap \Lang L(\Gamma_2)$ या भाषेतील
पूर्ण उपपत्ती~$\Gamma$ समाविष्ट असेल,
तर $\Gamma_1 \cup \Gamma_2$ सुसंगत असते.

पूर्ण उपपत्ती म्हणजे तिच्या भाषेतील
प्रत्येक !!{sentence}~$!A$ साठी
$\Gamma \Proves !A$ किंवा
$\Gamma \Proves \lnot !A$ सिद्ध होणारी उपपत्ती,
हे आठवा. $\Gamma_1$ आणि~$\Gamma_2$
या $\Gamma$ च्या सुसंगत विस्तारे
असल्याने $\Gamma$ सुसंगतच असली पाहिजे.
म्हणून $!A$ हे
$\Lang L(\Gamma_1) \cap \Lang
L(\Gamma_2)$ या भाषेतील
!!a{sentence} असेल, तर
$\Gamma_1 \Entails !A$ आणि
$\Gamma_2
\Entails \lnot !A$ दोन्ही शक्य नाहीत.
तसे एखादे~$!A$ आहे असे समजा.
$\Gamma_1 \Entails !A$ असेल, तर
$\Gamma \Entails !A$.
अन्यथा, $\Gamma$ पूर्ण असल्याने
$\Gamma \Entails \lnot !A$;
आणि $\Gamma \subseteq \Gamma_1$
असल्याने $\Gamma_1 \Entails \lnot !A$,
म्हणजे $\Gamma_1$ विसंगत ठरेल.
म्हणून $\Gamma \Entails !A$;
आता $\Gamma
\subseteq \Gamma_2$ असल्याने
$\Gamma_2 \Entails !A$.
त्यामुळे $\Gamma \Entails/
\lnot !A$; अन्यथा $\Gamma_2$
विसंगत ठरेल.

संयुक्त !!{language}~
$\Lang L(\Gamma_1) \cap \Lang L(\Gamma_2)$
मधील असे !!a{sentence}~$!A$
नसणे की $\Gamma_1
\Entails !A$ आणि
$\Gamma_2 \Entails \lnot !A$,
हेच सिद्धतेसाठी पुरेसे आहे.
आता माएहाराचे उपप्रमेय वापरून
ही सिद्धता सिद्धता-उपपत्तीय
पद्धतीने देऊ.

\begin{thm}[रॉबिन्सनचे संयुक्त सुसंगतता प्रमेय]
$\Gamma_1$ आणि $\Gamma_2$
या सुसंगत उपपत्ती असू द्या.
$\Lang L(\Gamma_1) \cap \Lang L(\Gamma_2)$
या !!{language} मधील कोणत्याही~$!A$
साठी $\Gamma_1 \Entails !A$
आणि $\Gamma_2 \Entails \lnot !A$
दोन्ही होत नसतील, तर
$\Gamma_1 \cup \Gamma_2$ सुसंगत असते.
\end{thm}

\begin{proof}
  उलट, $\Gamma_1 \cup \Gamma_2$
  विसंगत आहे असे समजा.
  मग $\Gamma_1, \Gamma_2 \Sequent \ $
  या क्रमवर्तीची !!a{proof} आहे.
  $\maeh{\Gamma_1}{\Gamma_2}
  \Sequent \maeh{\ }{\ }$ याला
  माएहाराचे उपप्रमेय लावून, आवश्यक असल्यास वर
  सांगितल्याप्रमाणे अंतर्वेशकातील मुक्त चरे बांधल्याने,
  $\Lang L(\Gamma_1) \cap \Lang L(\Gamma_2)$
  या !!{language} मधील असे
  !!a{sentence}~$!A$ मिळते की
  $\Proves \Gamma_1 \Sequent !A$
  आणि $\Proves !A, \Gamma_2 \Sequent
  \quad$. दुसऱ्या क्रमवर्तीवरून
  $\Proves \Gamma_2 \Sequent \lnot !A$
  मिळते. पण असे !!{sentence}
  अस्तित्वात नाही, हेच आपण
  गृहीत धरले होते.
\end{proof}

वर आपण उपपत्ती $\Gamma$,
$\Gamma_1$ आणि $\Gamma_2$
सांत आहेत असे नकळत गृहीत धरले.
संहततेमुळे हे निष्कर्ष
अनंत उपपत्तींनाही लागू होतात.

\end{document}
